% ==========================================================================
%  OhmTools Lernhelfer – Vorlage für Übungsblätter im Klausurstil
%  ------------------------------------------------------------------------
%  Zwei Versionen aus EINER Datei:
%
%   Angabenblatt:    pdflatex uebungsblatt.tex          (3x ausführen)
%   Lösungsversion:  pdflatex -jobname=uebungsblatt-loesung "\def\mitloesung{}% ==========================================================================
%  OhmTools Lernhelfer – Vorlage für Übungsblätter im Klausurstil
%  ------------------------------------------------------------------------
%  Zwei Versionen aus EINER Datei:
%
%   Angabenblatt:    pdflatex uebungsblatt.tex          (3x ausführen)
%   Lösungsversion:  pdflatex -jobname=uebungsblatt-loesung "\def\mitloesung{}% ==========================================================================
%  OhmTools Lernhelfer – Vorlage für Übungsblätter im Klausurstil
%  ------------------------------------------------------------------------
%  Zwei Versionen aus EINER Datei:
%
%   Angabenblatt:    pdflatex uebungsblatt.tex          (3x ausführen)
%   Lösungsversion:  pdflatex -jobname=uebungsblatt-loesung "\def\mitloesung{}% ==========================================================================
%  OhmTools Lernhelfer – Vorlage für Übungsblätter im Klausurstil
%  ------------------------------------------------------------------------
%  Zwei Versionen aus EINER Datei:
%
%   Angabenblatt:    pdflatex uebungsblatt.tex          (3x ausführen)
%   Lösungsversion:  pdflatex -jobname=uebungsblatt-loesung "\def\mitloesung{}\input{uebungsblatt.tex}"
%                    (ebenfalls 3x ausführen)   – oder einfach: sh build.sh
%
%  Alternativ unten die Zeile  % \printanswers  einkommentieren.
%
%  Benötigt: TeX Live mit exam, circuitikz, pgfplots, siunitx, tcolorbox,
%            listings, tex-gyre (Pagella/Heros), lmodern (Code-Schrift),
%            babel-german, hyperref
% ==========================================================================
\documentclass[11pt,a4paper,addpoints]{exam}

% --- Lösungen ein/aus ----------------------------------------------------
% \printanswers
\ifdefined\mitloesung\printanswers\fi

% ==========================================================================
%  KOPFDATEN – hier anpassen
% ==========================================================================
\newcommand{\modulname}{Ingenieurgrundlagen}
\newcommand{\dauer}{90 Minuten}
\newcommand{\hilfsmittel}{nicht programmierbarer Taschenrechner, 1 Formelblatt (A4)}

% ==========================================================================
%  PAKETE
% ==========================================================================
\usepackage[T1]{fontenc}
\usepackage[ngerman]{babel}
\usepackage[sc]{mathpazo}            % Mathe-Schrift passend zu Palatino
\usepackage{tgpagella}               % Fließtext: TeX Gyre Pagella
\usepackage[scale=0.92]{tgheros}     % Überschriften: TeX Gyre Heros
\renewcommand{\ttdefault}{lmtt}      % Code: Latin Modern Mono
\usepackage{microtype}
\usepackage{amsmath,amssymb}
\usepackage{xcolor}
\usepackage[most]{tcolorbox}
\usepackage{siunitx}
\sisetup{locale=DE, per-mode=symbol}
\usepackage{tikz}
\usetikzlibrary{arrows.meta,calc}
\usepackage{pgfplots}
\pgfplotsset{compat=1.18}
\usepackage[european]{circuitikz}
\usepackage{listings}
\usepackage[hidelinks]{hyperref}
\hypersetup{pdftitle={OhmTools Lernhelfer Übungsblatt für \modulname}}

% --- Farben ---------------------------------------------------------------
\definecolor{akzent}{HTML}{1F4E79}   % Hauptfarbe (dunkelblau)
\definecolor{akzenthell}{HTML}{E8EEF5}
\definecolor{grau}{HTML}{6B7280}
\definecolor{loesung}{HTML}{1B7F3B}
\definecolor{GridColor}{HTML}{C9D3DE}

% ==========================================================================
%  LAYOUT DER EXAM-KLASSE
% ==========================================================================
\extrawidth{0.4in}
\extraheadheight{-0.1in}
\extrafootheight{0.1in}

% Kopf- und Fußzeilen
\pagestyle{headandfoot}
\firstpageheader{}{}{}
\runningheader{\sffamily\small\color{grau}OhmTools Lernhelfer}
              {\ifprintanswers\sffamily\small\bfseries\color{loesung}LÖSUNG\fi}
              {\sffamily\small\color{grau}Übungsblatt \modulname}
\runningheadrule
\footer{}{\sffamily\small\color{grau}Seite \thepage\ von \numpages}{}

% Aufgabenkopf: „Aufgabe 1  Titel ............ 12 Punkte“
\newcommand{\aufgabenkopf}{%
  \begin{minipage}{\textwidth}
    \vspace{3mm}
    {\sffamily\bfseries\large\color{akzent}Aufgabe \thequestion\quad\thequestiontitle}%
    \hfill{\sffamily\small\color{grau}\totalpoints\ \points}\par\vspace{1mm}
    {\color{akzent}\rule{\textwidth}{0.8pt}}\rule[-3mm]{0pt}{0pt}
  \end{minipage}}
\qformat{\aufgabenkopf}

% Jede Aufgabe beginnt auf einer neuen Seite, die erste unter dem Titelkopf.
% So lässt sich jede Aufgabe am Tablet ohne Scrollen bearbeiten, der Rest
% der Seite bleibt frei.
\newcommand{\aufgabe}[1]{\ifnum\value{question}>0 \newpage\fi\titledquestion{#1}}

% Punkte bei Teilaufgaben
\pointpoints{Punkt}{Punkte}
\pointname{\,P.}
\pointformat{\textcolor{grau}{\sffamily\footnotesize[\thepoints]}}
\renewcommand{\partlabel}{\textbf{\thepartno)}}
\renewcommand{\partshook}{\setlength{\itemsep}{4pt}}

% Lösungen (erscheinen nur in der Lösungsversion)
\shadedsolutions
\definecolor{SolutionColor}{HTML}{EDF7F0}
\renewcommand{\solutiontitle}{\noindent\textbf{\sffamily\color{loesung}Lösung:}\enspace}

% Multiple Choice
\checkboxchar{$\square$}
\checkedchar{$\blacksquare$}
\CorrectChoiceEmphasis{\bfseries\color{loesung}}

% Code-Darstellung
\lstset{
  basicstyle=\ttfamily\small,
  columns=flexible, keepspaces=true,
  keywordstyle=\color{akzent}\bfseries,
  commentstyle=\color{grau}\itshape,
  stringstyle=\color{loesung},
  numbers=left, numberstyle=\tiny\color{grau}, numbersep=8pt,
  frame=leftline, rulecolor=\color{akzent}, framerule=1.2pt,
  xleftmargin=2.5em, aboveskip=6pt, belowskip=6pt,
  showstringspaces=false, tabsize=4,
  backgroundcolor=\color{akzenthell!60},
}

% Kasten für „Gegeben“
\newtcolorbox{gegeben}{enhanced, colback=akzenthell, colframe=akzenthell,
  boxrule=0pt, arc=2pt, left=3mm, right=3mm, top=1.5mm, bottom=1.5mm,
  fontupper=\small}

\begin{document}

% ==========================================================================
%  TITELKOPF
% ==========================================================================
\begin{tcolorbox}[enhanced, colback=akzent, colframe=akzent, coltext=white,
    arc=3pt, boxrule=0pt, left=6mm, right=6mm, top=4mm, bottom=4mm]
  {\sffamily\small OhmTools Lernhelfer\par}
  \vspace{2mm}
  {\sffamily\bfseries\LARGE Übungsblatt für \modulname\par}
  \ifprintanswers
    \vspace{2mm}
    {\sffamily\small\hfill\colorbox{white}{\sffamily\bfseries\color{loesung}\,LÖSUNGSVERSION\,}\par}
  \fi
\end{tcolorbox}

\vspace{1mm}
\noindent{\sffamily\small
  {\color{grau}Bearbeitungszeit:}~\dauer\hfill
  {\color{grau}Hilfsmittel:}~\hilfsmittel\hfill
  {\color{grau}Punkte:}~\numpoints\par}
\vspace{2mm}

% ==========================================================================
%  AUFGABEN
% ==========================================================================
\begin{questions}

% --------------------------------------------------------------------------
%  Aufgabe 1 – Mathematik
% --------------------------------------------------------------------------
\aufgabe{Kurvendiskussion}
Gegeben ist die Funktion
\[
  f(x) = x^3 - 6x^2 + 9x, \qquad x \in \mathbb{R}.
\]

\begin{parts}
\part[2] Bestimmen Sie die Nullstellen von $f$.
\begin{solution}
  $f(x) = x\,(x^2 - 6x + 9) = x\,(x-3)^2$.
  Nullstellen: $x_1 = 0$ (einfach), $x_2 = 3$ (doppelt).
\end{solution}

\part[4] Berechnen Sie die Extrempunkte und geben Sie jeweils deren Art an.
\begin{solution}
  $f'(x) = 3x^2 - 12x + 9 = 3\,(x-1)(x-3) \overset{!}{=} 0
   \;\Rightarrow\; x = 1$ oder $x = 3$.\\
  $f''(x) = 6x - 12$: \quad $f''(1) = -6 < 0 \Rightarrow$ Hochpunkt $H(1\,|\,4)$,
  \quad $f''(3) = 6 > 0 \Rightarrow$ Tiefpunkt $T(3\,|\,0)$.
\end{solution}

\part[2] Bestimmen Sie den Wendepunkt.
\begin{solution}
  $f''(x) = 0 \Rightarrow x = 2$, \ $f'''(2) = 6 \neq 0$.
  Wendepunkt $W(2\,|\,2)$.
\end{solution}

\part[3] Berechnen Sie den Inhalt der Fläche, die der Graph von $f$ im Intervall
$[0;\,3]$ mit der $x$-Achse einschließt.
\begin{solution}
  Da $f(x) = x(x-3)^2 \ge 0$ auf $[0;3]$:
  \[
    A = \int_0^3 f(x)\,\mathrm{d}x
      = \Big[\tfrac{x^4}{4} - 2x^3 + \tfrac{9}{2}x^2\Big]_0^3
      = \tfrac{81}{4} - 54 + \tfrac{81}{2} = \tfrac{27}{4} = 6{,}75 \text{ FE}.
  \]
\end{solution}

\part[2] Skizzieren Sie den Graphen von $f$ im folgenden Koordinatensystem
und markieren Sie die berechneten Punkte.

\begin{center}
\begin{tikzpicture}
  \begin{axis}[
      width=10cm, height=5.5cm,
      axis lines=middle, xmin=-0.6, xmax=4.6, ymin=-1.6, ymax=5.6,
      xtick={-0.5,0,...,4.5}, ytick={-1,...,5},
      xticklabels={,0,,1,,2,,3,,4,},
      grid=both, grid style={GridColor, line width=0.2pt},
      minor tick num=1, minor grid style={GridColor!50},
      xlabel={$x$}, ylabel={$y$},
      tick label style={font=\footnotesize},
      clip=true]
    \ifprintanswers
      \addplot[akzent, very thick, domain=-0.3:4.25, samples=120] {x^3 - 6*x^2 + 9*x};
      \addplot[only marks, mark=*, mark size=2pt, loesung]
        coordinates {(0,0) (3,0) (1,4) (2,2)};
      \node[loesung, above right, font=\footnotesize] at (axis cs:1,4) {$H$};
      \node[loesung, above right, font=\footnotesize] at (axis cs:2,2) {$W$};
      \node[loesung, below right, font=\footnotesize] at (axis cs:3,0) {$T$};
    \fi
  \end{axis}
\end{tikzpicture}
\end{center}
\end{parts}

% --------------------------------------------------------------------------
%  Aufgabe 2 – Elektrotechnik
% --------------------------------------------------------------------------
\aufgabe{Gleichstromnetzwerk}
Gegeben ist die folgende Schaltung.

\begin{center}
\begin{circuitikz}[line width=0.7pt]
  \draw (0,0) to[V, v<=$U_0$] (0,3)
        to[R, l=$R_1$, i>^=$I$] (3.5,3)
        to[short, -, i>^=$I_3$] (6.5,3)
        to[R, l_=$R_3$, v^=$U_P$] (6.5,0)
        -- (0,0);
  \draw (3.5,3) to[R, l=$R_2$, i>_=$I_2$] (3.5,0);
  \fill (3.5,3) circle (2pt) (3.5,0) circle (2pt);
\end{circuitikz}
\end{center}

\begin{gegeben}
  \textbf{Gegeben:}\quad
  $U_0 = \qty{12}{\volt}$, \quad
  $R_1 = \qty{100}{\ohm}$, \quad
  $R_2 = \qty{300}{\ohm}$, \quad
  $R_3 = \qty{150}{\ohm}$
\end{gegeben}

\begin{parts}
\part[3] Berechnen Sie den Gesamtwiderstand $R_\text{ges}$ der Schaltung.
\begin{solution}
  $R_{23} = \dfrac{R_2 R_3}{R_2 + R_3} = \dfrac{300 \cdot 150}{450}\,\unit{\ohm} = \qty{100}{\ohm}$,
  \quad $R_\text{ges} = R_1 + R_{23} = \qty{200}{\ohm}$.
\end{solution}

\part[2] Berechnen Sie den Gesamtstrom $I$.
\begin{solution}
  $I = \dfrac{U_0}{R_\text{ges}} = \dfrac{\qty{12}{\volt}}{\qty{200}{\ohm}} = \qty{60}{\milli\ampere}$.
\end{solution}

\part[4] Bestimmen Sie die Spannung $U_P$ an der Parallelschaltung sowie die
Teilströme $I_2$ und $I_3$.
\begin{solution}
  $U_P = I \cdot R_{23} = \qty{60}{\milli\ampere} \cdot \qty{100}{\ohm} = \qty{6}{\volt}$.\\
  $I_2 = U_P / R_2 = \qty{20}{\milli\ampere}$, \quad
  $I_3 = U_P / R_3 = \qty{40}{\milli\ampere}$
  \quad (Probe: $I_2 + I_3 = I$ \checkmark).
\end{solution}

\part[2] Welche Leistung wird im Widerstand $R_1$ umgesetzt?
\begin{solution}
  $P_1 = I^2 R_1 = (\qty{0,06}{\ampere})^2 \cdot \qty{100}{\ohm} = \qty{0,36}{\watt}$.
\end{solution}
\end{parts}

% --------------------------------------------------------------------------
%  Aufgabe 3 – Physik
% --------------------------------------------------------------------------
\aufgabe{Schräger Wurf}
Ein Ball wird vom Boden aus mit der Anfangsgeschwindigkeit $v_0$ unter dem Winkel
$\alpha$ gegen die Horizontale abgeworfen. Der Luftwiderstand wird vernachlässigt.

\begin{center}
\begin{tikzpicture}[>=Stealth]
  \draw[->, grau] (-0.3,0) -- (9.8,0) node[below] {$x$};
  \draw[->, grau] (0,-0.3) -- (0,3.4) node[left] {$y$};
  \draw[akzent, very thick] (0,0) parabola bend (4.5,2.6) (9,0);
  \draw[->, very thick] (0,0) -- (49:1.9) node[above left] {$\vec v_0$};
  \draw (1,0) arc (0:49:1);
  \node at (24:1.3) {$\alpha$};
  \draw[<->, dashed, grau] (4.5,0) -- node[right, black] {$h_\text{max}$} (4.5,2.6);
  \draw[<->] (0,-0.6) -- node[fill=white, inner sep=2pt] {$w$} (9,-0.6);
  \fill[akzent] (4.5,2.6) circle (2pt);
  \node[grau, font=\footnotesize, anchor=north east] at (9.8,3.4) {Skizze nicht maßstäblich};
\end{tikzpicture}
\end{center}

\begin{gegeben}
  \textbf{Gegeben:}\quad
  $v_0 = \qty{20}{\metre\per\second}$, \quad
  $\alpha = \ang{30}$, \quad
  $g = \qty{9,81}{\metre\per\second\squared}$
\end{gegeben}

\begin{parts}
\part[2] Nach welcher Zeit $t_s$ erreicht der Ball seine maximale Höhe?
\begin{solution}
  Am Scheitel gilt $v_y = 0$:\quad
  $t_s = \dfrac{v_0 \sin\alpha}{g} = \dfrac{\qty{10}{\metre\per\second}}{\qty{9,81}{\metre\per\second\squared}} \approx \qty{1,02}{\second}$.
\end{solution}

\part[3] Berechnen Sie die maximale Höhe $h_\text{max}$.
\begin{solution}
  $h_\text{max} = \dfrac{v_0^2 \sin^2\alpha}{2g}
   = \dfrac{(\qty{10}{\metre\per\second})^2}{2 \cdot \qty{9,81}{\metre\per\second\squared}}
   \approx \qty{5,10}{\metre}$.
\end{solution}

\part[3] Berechnen Sie die Wurfweite $w$.
\begin{solution}
  $w = \dfrac{v_0^2 \sin(2\alpha)}{g}
   = \dfrac{\qty{400}{\metre\squared\per\second\squared} \cdot \sin\ang{60}}{\qty{9,81}{\metre\per\second\squared}}
   \approx \qty{35,3}{\metre}$.
\end{solution}

\part[2] Unter welchem Abwurfwinkel wird die Wurfweite maximal? Begründen Sie kurz.
\begin{solution}
  $w \sim \sin(2\alpha)$ wird maximal für $2\alpha = \ang{90}$, also $\alpha = \ang{45}$.
\end{solution}
\end{parts}

% --------------------------------------------------------------------------
%  Aufgabe 4 – Informatik
% --------------------------------------------------------------------------
\aufgabe{Rekursion und Zahlensysteme}
Betrachten Sie die folgende Python-Funktion:

\begin{lstlisting}[language=Python]
def f(n):
    if n <= 1:
        return n
    return f(n - 1) + f(n - 2)
\end{lstlisting}

\begin{parts}
\part[2] Welchen Wert liefert \lstinline{f(6)}? Geben Sie die Zwischenwerte an.
\begin{solution}
  $f(0..6) = 0,\,1,\,1,\,2,\,3,\,5,\,8$, \ also \lstinline{f(6)} $= 8$.
\end{solution}

\part[1] Was berechnet die Funktion allgemein?
\begin{solution}
  Die $n$-te Fibonacci-Zahl.
\end{solution}

\part[3] Geben Sie die Laufzeitkomplexität in $\mathcal{O}$-Notation an und
begründen Sie. Wie lässt sich die Laufzeit verbessern?
\begin{solution}
  $\mathcal{O}(2^n)$: Jeder Aufruf erzeugt bis zu zwei weitere Aufrufe, der
  Rekursionsbaum hat die Tiefe $n$, und Teilergebnisse werden mehrfach berechnet.
  Mit Memoisation oder einer iterativen Lösung sinkt die Laufzeit auf $\mathcal{O}(n)$.
\end{solution}

\part[1] Welche Datenstruktur arbeitet nach dem LIFO-Prinzip? Kreuzen Sie an.
\par\smallskip
\begin{oneparcheckboxes}
  \choice Warteschlange (Queue)
  \CorrectChoice Stapel (Stack)
  \choice Heap
  \choice Hashtabelle
\end{oneparcheckboxes}

\part[2] Wandeln Sie die Dezimalzahl $45$ in das Binär- und das Hexadezimalsystem um.
\begin{solution}
  $45 = 32 + 8 + 4 + 1 = 101101_2$, \quad $0010\,1101_2 = 2\mathrm{D}_{16}$.
\end{solution}
\end{parts}

\end{questions}

\end{document}
"
%                    (ebenfalls 3x ausführen)   – oder einfach: sh build.sh
%
%  Alternativ unten die Zeile  % \printanswers  einkommentieren.
%
%  Benötigt: TeX Live mit exam, circuitikz, pgfplots, siunitx, tcolorbox,
%            listings, tex-gyre (Pagella/Heros), lmodern (Code-Schrift),
%            babel-german, hyperref
% ==========================================================================
\documentclass[11pt,a4paper,addpoints]{exam}

% --- Lösungen ein/aus ----------------------------------------------------
% \printanswers
\ifdefined\mitloesung\printanswers\fi

% ==========================================================================
%  KOPFDATEN – hier anpassen
% ==========================================================================
\newcommand{\modulname}{Ingenieurgrundlagen}
\newcommand{\dauer}{90 Minuten}
\newcommand{\hilfsmittel}{nicht programmierbarer Taschenrechner, 1 Formelblatt (A4)}

% ==========================================================================
%  PAKETE
% ==========================================================================
\usepackage[T1]{fontenc}
\usepackage[ngerman]{babel}
\usepackage[sc]{mathpazo}            % Mathe-Schrift passend zu Palatino
\usepackage{tgpagella}               % Fließtext: TeX Gyre Pagella
\usepackage[scale=0.92]{tgheros}     % Überschriften: TeX Gyre Heros
\renewcommand{\ttdefault}{lmtt}      % Code: Latin Modern Mono
\usepackage{microtype}
\usepackage{amsmath,amssymb}
\usepackage{xcolor}
\usepackage[most]{tcolorbox}
\usepackage{siunitx}
\sisetup{locale=DE, per-mode=symbol}
\usepackage{tikz}
\usetikzlibrary{arrows.meta,calc}
\usepackage{pgfplots}
\pgfplotsset{compat=1.18}
\usepackage[european]{circuitikz}
\usepackage{listings}
\usepackage[hidelinks]{hyperref}
\hypersetup{pdftitle={OhmTools Lernhelfer Übungsblatt für \modulname}}

% --- Farben ---------------------------------------------------------------
\definecolor{akzent}{HTML}{1F4E79}   % Hauptfarbe (dunkelblau)
\definecolor{akzenthell}{HTML}{E8EEF5}
\definecolor{grau}{HTML}{6B7280}
\definecolor{loesung}{HTML}{1B7F3B}
\definecolor{GridColor}{HTML}{C9D3DE}

% ==========================================================================
%  LAYOUT DER EXAM-KLASSE
% ==========================================================================
\extrawidth{0.4in}
\extraheadheight{-0.1in}
\extrafootheight{0.1in}

% Kopf- und Fußzeilen
\pagestyle{headandfoot}
\firstpageheader{}{}{}
\runningheader{\sffamily\small\color{grau}OhmTools Lernhelfer}
              {\ifprintanswers\sffamily\small\bfseries\color{loesung}LÖSUNG\fi}
              {\sffamily\small\color{grau}Übungsblatt \modulname}
\runningheadrule
\footer{}{\sffamily\small\color{grau}Seite \thepage\ von \numpages}{}

% Aufgabenkopf: „Aufgabe 1  Titel ............ 12 Punkte“
\newcommand{\aufgabenkopf}{%
  \begin{minipage}{\textwidth}
    \vspace{3mm}
    {\sffamily\bfseries\large\color{akzent}Aufgabe \thequestion\quad\thequestiontitle}%
    \hfill{\sffamily\small\color{grau}\totalpoints\ \points}\par\vspace{1mm}
    {\color{akzent}\rule{\textwidth}{0.8pt}}\rule[-3mm]{0pt}{0pt}
  \end{minipage}}
\qformat{\aufgabenkopf}

% Jede Aufgabe beginnt auf einer neuen Seite, die erste unter dem Titelkopf.
% So lässt sich jede Aufgabe am Tablet ohne Scrollen bearbeiten, der Rest
% der Seite bleibt frei.
\newcommand{\aufgabe}[1]{\ifnum\value{question}>0 \newpage\fi\titledquestion{#1}}

% Punkte bei Teilaufgaben
\pointpoints{Punkt}{Punkte}
\pointname{\,P.}
\pointformat{\textcolor{grau}{\sffamily\footnotesize[\thepoints]}}
\renewcommand{\partlabel}{\textbf{\thepartno)}}
\renewcommand{\partshook}{\setlength{\itemsep}{4pt}}

% Lösungen (erscheinen nur in der Lösungsversion)
\shadedsolutions
\definecolor{SolutionColor}{HTML}{EDF7F0}
\renewcommand{\solutiontitle}{\noindent\textbf{\sffamily\color{loesung}Lösung:}\enspace}

% Multiple Choice
\checkboxchar{$\square$}
\checkedchar{$\blacksquare$}
\CorrectChoiceEmphasis{\bfseries\color{loesung}}

% Code-Darstellung
\lstset{
  basicstyle=\ttfamily\small,
  columns=flexible, keepspaces=true,
  keywordstyle=\color{akzent}\bfseries,
  commentstyle=\color{grau}\itshape,
  stringstyle=\color{loesung},
  numbers=left, numberstyle=\tiny\color{grau}, numbersep=8pt,
  frame=leftline, rulecolor=\color{akzent}, framerule=1.2pt,
  xleftmargin=2.5em, aboveskip=6pt, belowskip=6pt,
  showstringspaces=false, tabsize=4,
  backgroundcolor=\color{akzenthell!60},
}

% Kasten für „Gegeben“
\newtcolorbox{gegeben}{enhanced, colback=akzenthell, colframe=akzenthell,
  boxrule=0pt, arc=2pt, left=3mm, right=3mm, top=1.5mm, bottom=1.5mm,
  fontupper=\small}

\begin{document}

% ==========================================================================
%  TITELKOPF
% ==========================================================================
\begin{tcolorbox}[enhanced, colback=akzent, colframe=akzent, coltext=white,
    arc=3pt, boxrule=0pt, left=6mm, right=6mm, top=4mm, bottom=4mm]
  {\sffamily\small OhmTools Lernhelfer\par}
  \vspace{2mm}
  {\sffamily\bfseries\LARGE Übungsblatt für \modulname\par}
  \ifprintanswers
    \vspace{2mm}
    {\sffamily\small\hfill\colorbox{white}{\sffamily\bfseries\color{loesung}\,LÖSUNGSVERSION\,}\par}
  \fi
\end{tcolorbox}

\vspace{1mm}
\noindent{\sffamily\small
  {\color{grau}Bearbeitungszeit:}~\dauer\hfill
  {\color{grau}Hilfsmittel:}~\hilfsmittel\hfill
  {\color{grau}Punkte:}~\numpoints\par}
\vspace{2mm}

% ==========================================================================
%  AUFGABEN
% ==========================================================================
\begin{questions}

% --------------------------------------------------------------------------
%  Aufgabe 1 – Mathematik
% --------------------------------------------------------------------------
\aufgabe{Kurvendiskussion}
Gegeben ist die Funktion
\[
  f(x) = x^3 - 6x^2 + 9x, \qquad x \in \mathbb{R}.
\]

\begin{parts}
\part[2] Bestimmen Sie die Nullstellen von $f$.
\begin{solution}
  $f(x) = x\,(x^2 - 6x + 9) = x\,(x-3)^2$.
  Nullstellen: $x_1 = 0$ (einfach), $x_2 = 3$ (doppelt).
\end{solution}

\part[4] Berechnen Sie die Extrempunkte und geben Sie jeweils deren Art an.
\begin{solution}
  $f'(x) = 3x^2 - 12x + 9 = 3\,(x-1)(x-3) \overset{!}{=} 0
   \;\Rightarrow\; x = 1$ oder $x = 3$.\\
  $f''(x) = 6x - 12$: \quad $f''(1) = -6 < 0 \Rightarrow$ Hochpunkt $H(1\,|\,4)$,
  \quad $f''(3) = 6 > 0 \Rightarrow$ Tiefpunkt $T(3\,|\,0)$.
\end{solution}

\part[2] Bestimmen Sie den Wendepunkt.
\begin{solution}
  $f''(x) = 0 \Rightarrow x = 2$, \ $f'''(2) = 6 \neq 0$.
  Wendepunkt $W(2\,|\,2)$.
\end{solution}

\part[3] Berechnen Sie den Inhalt der Fläche, die der Graph von $f$ im Intervall
$[0;\,3]$ mit der $x$-Achse einschließt.
\begin{solution}
  Da $f(x) = x(x-3)^2 \ge 0$ auf $[0;3]$:
  \[
    A = \int_0^3 f(x)\,\mathrm{d}x
      = \Big[\tfrac{x^4}{4} - 2x^3 + \tfrac{9}{2}x^2\Big]_0^3
      = \tfrac{81}{4} - 54 + \tfrac{81}{2} = \tfrac{27}{4} = 6{,}75 \text{ FE}.
  \]
\end{solution}

\part[2] Skizzieren Sie den Graphen von $f$ im folgenden Koordinatensystem
und markieren Sie die berechneten Punkte.

\begin{center}
\begin{tikzpicture}
  \begin{axis}[
      width=10cm, height=5.5cm,
      axis lines=middle, xmin=-0.6, xmax=4.6, ymin=-1.6, ymax=5.6,
      xtick={-0.5,0,...,4.5}, ytick={-1,...,5},
      xticklabels={,0,,1,,2,,3,,4,},
      grid=both, grid style={GridColor, line width=0.2pt},
      minor tick num=1, minor grid style={GridColor!50},
      xlabel={$x$}, ylabel={$y$},
      tick label style={font=\footnotesize},
      clip=true]
    \ifprintanswers
      \addplot[akzent, very thick, domain=-0.3:4.25, samples=120] {x^3 - 6*x^2 + 9*x};
      \addplot[only marks, mark=*, mark size=2pt, loesung]
        coordinates {(0,0) (3,0) (1,4) (2,2)};
      \node[loesung, above right, font=\footnotesize] at (axis cs:1,4) {$H$};
      \node[loesung, above right, font=\footnotesize] at (axis cs:2,2) {$W$};
      \node[loesung, below right, font=\footnotesize] at (axis cs:3,0) {$T$};
    \fi
  \end{axis}
\end{tikzpicture}
\end{center}
\end{parts}

% --------------------------------------------------------------------------
%  Aufgabe 2 – Elektrotechnik
% --------------------------------------------------------------------------
\aufgabe{Gleichstromnetzwerk}
Gegeben ist die folgende Schaltung.

\begin{center}
\begin{circuitikz}[line width=0.7pt]
  \draw (0,0) to[V, v<=$U_0$] (0,3)
        to[R, l=$R_1$, i>^=$I$] (3.5,3)
        to[short, -, i>^=$I_3$] (6.5,3)
        to[R, l_=$R_3$, v^=$U_P$] (6.5,0)
        -- (0,0);
  \draw (3.5,3) to[R, l=$R_2$, i>_=$I_2$] (3.5,0);
  \fill (3.5,3) circle (2pt) (3.5,0) circle (2pt);
\end{circuitikz}
\end{center}

\begin{gegeben}
  \textbf{Gegeben:}\quad
  $U_0 = \qty{12}{\volt}$, \quad
  $R_1 = \qty{100}{\ohm}$, \quad
  $R_2 = \qty{300}{\ohm}$, \quad
  $R_3 = \qty{150}{\ohm}$
\end{gegeben}

\begin{parts}
\part[3] Berechnen Sie den Gesamtwiderstand $R_\text{ges}$ der Schaltung.
\begin{solution}
  $R_{23} = \dfrac{R_2 R_3}{R_2 + R_3} = \dfrac{300 \cdot 150}{450}\,\unit{\ohm} = \qty{100}{\ohm}$,
  \quad $R_\text{ges} = R_1 + R_{23} = \qty{200}{\ohm}$.
\end{solution}

\part[2] Berechnen Sie den Gesamtstrom $I$.
\begin{solution}
  $I = \dfrac{U_0}{R_\text{ges}} = \dfrac{\qty{12}{\volt}}{\qty{200}{\ohm}} = \qty{60}{\milli\ampere}$.
\end{solution}

\part[4] Bestimmen Sie die Spannung $U_P$ an der Parallelschaltung sowie die
Teilströme $I_2$ und $I_3$.
\begin{solution}
  $U_P = I \cdot R_{23} = \qty{60}{\milli\ampere} \cdot \qty{100}{\ohm} = \qty{6}{\volt}$.\\
  $I_2 = U_P / R_2 = \qty{20}{\milli\ampere}$, \quad
  $I_3 = U_P / R_3 = \qty{40}{\milli\ampere}$
  \quad (Probe: $I_2 + I_3 = I$ \checkmark).
\end{solution}

\part[2] Welche Leistung wird im Widerstand $R_1$ umgesetzt?
\begin{solution}
  $P_1 = I^2 R_1 = (\qty{0,06}{\ampere})^2 \cdot \qty{100}{\ohm} = \qty{0,36}{\watt}$.
\end{solution}
\end{parts}

% --------------------------------------------------------------------------
%  Aufgabe 3 – Physik
% --------------------------------------------------------------------------
\aufgabe{Schräger Wurf}
Ein Ball wird vom Boden aus mit der Anfangsgeschwindigkeit $v_0$ unter dem Winkel
$\alpha$ gegen die Horizontale abgeworfen. Der Luftwiderstand wird vernachlässigt.

\begin{center}
\begin{tikzpicture}[>=Stealth]
  \draw[->, grau] (-0.3,0) -- (9.8,0) node[below] {$x$};
  \draw[->, grau] (0,-0.3) -- (0,3.4) node[left] {$y$};
  \draw[akzent, very thick] (0,0) parabola bend (4.5,2.6) (9,0);
  \draw[->, very thick] (0,0) -- (49:1.9) node[above left] {$\vec v_0$};
  \draw (1,0) arc (0:49:1);
  \node at (24:1.3) {$\alpha$};
  \draw[<->, dashed, grau] (4.5,0) -- node[right, black] {$h_\text{max}$} (4.5,2.6);
  \draw[<->] (0,-0.6) -- node[fill=white, inner sep=2pt] {$w$} (9,-0.6);
  \fill[akzent] (4.5,2.6) circle (2pt);
  \node[grau, font=\footnotesize, anchor=north east] at (9.8,3.4) {Skizze nicht maßstäblich};
\end{tikzpicture}
\end{center}

\begin{gegeben}
  \textbf{Gegeben:}\quad
  $v_0 = \qty{20}{\metre\per\second}$, \quad
  $\alpha = \ang{30}$, \quad
  $g = \qty{9,81}{\metre\per\second\squared}$
\end{gegeben}

\begin{parts}
\part[2] Nach welcher Zeit $t_s$ erreicht der Ball seine maximale Höhe?
\begin{solution}
  Am Scheitel gilt $v_y = 0$:\quad
  $t_s = \dfrac{v_0 \sin\alpha}{g} = \dfrac{\qty{10}{\metre\per\second}}{\qty{9,81}{\metre\per\second\squared}} \approx \qty{1,02}{\second}$.
\end{solution}

\part[3] Berechnen Sie die maximale Höhe $h_\text{max}$.
\begin{solution}
  $h_\text{max} = \dfrac{v_0^2 \sin^2\alpha}{2g}
   = \dfrac{(\qty{10}{\metre\per\second})^2}{2 \cdot \qty{9,81}{\metre\per\second\squared}}
   \approx \qty{5,10}{\metre}$.
\end{solution}

\part[3] Berechnen Sie die Wurfweite $w$.
\begin{solution}
  $w = \dfrac{v_0^2 \sin(2\alpha)}{g}
   = \dfrac{\qty{400}{\metre\squared\per\second\squared} \cdot \sin\ang{60}}{\qty{9,81}{\metre\per\second\squared}}
   \approx \qty{35,3}{\metre}$.
\end{solution}

\part[2] Unter welchem Abwurfwinkel wird die Wurfweite maximal? Begründen Sie kurz.
\begin{solution}
  $w \sim \sin(2\alpha)$ wird maximal für $2\alpha = \ang{90}$, also $\alpha = \ang{45}$.
\end{solution}
\end{parts}

% --------------------------------------------------------------------------
%  Aufgabe 4 – Informatik
% --------------------------------------------------------------------------
\aufgabe{Rekursion und Zahlensysteme}
Betrachten Sie die folgende Python-Funktion:

\begin{lstlisting}[language=Python]
def f(n):
    if n <= 1:
        return n
    return f(n - 1) + f(n - 2)
\end{lstlisting}

\begin{parts}
\part[2] Welchen Wert liefert \lstinline{f(6)}? Geben Sie die Zwischenwerte an.
\begin{solution}
  $f(0..6) = 0,\,1,\,1,\,2,\,3,\,5,\,8$, \ also \lstinline{f(6)} $= 8$.
\end{solution}

\part[1] Was berechnet die Funktion allgemein?
\begin{solution}
  Die $n$-te Fibonacci-Zahl.
\end{solution}

\part[3] Geben Sie die Laufzeitkomplexität in $\mathcal{O}$-Notation an und
begründen Sie. Wie lässt sich die Laufzeit verbessern?
\begin{solution}
  $\mathcal{O}(2^n)$: Jeder Aufruf erzeugt bis zu zwei weitere Aufrufe, der
  Rekursionsbaum hat die Tiefe $n$, und Teilergebnisse werden mehrfach berechnet.
  Mit Memoisation oder einer iterativen Lösung sinkt die Laufzeit auf $\mathcal{O}(n)$.
\end{solution}

\part[1] Welche Datenstruktur arbeitet nach dem LIFO-Prinzip? Kreuzen Sie an.
\par\smallskip
\begin{oneparcheckboxes}
  \choice Warteschlange (Queue)
  \CorrectChoice Stapel (Stack)
  \choice Heap
  \choice Hashtabelle
\end{oneparcheckboxes}

\part[2] Wandeln Sie die Dezimalzahl $45$ in das Binär- und das Hexadezimalsystem um.
\begin{solution}
  $45 = 32 + 8 + 4 + 1 = 101101_2$, \quad $0010\,1101_2 = 2\mathrm{D}_{16}$.
\end{solution}
\end{parts}

\end{questions}

\end{document}
"
%                    (ebenfalls 3x ausführen)   – oder einfach: sh build.sh
%
%  Alternativ unten die Zeile  % \printanswers  einkommentieren.
%
%  Benötigt: TeX Live mit exam, circuitikz, pgfplots, siunitx, tcolorbox,
%            listings, tex-gyre (Pagella/Heros), lmodern (Code-Schrift),
%            babel-german, hyperref
% ==========================================================================
\documentclass[11pt,a4paper,addpoints]{exam}

% --- Lösungen ein/aus ----------------------------------------------------
% \printanswers
\ifdefined\mitloesung\printanswers\fi

% ==========================================================================
%  KOPFDATEN – hier anpassen
% ==========================================================================
\newcommand{\modulname}{Ingenieurgrundlagen}
\newcommand{\dauer}{90 Minuten}
\newcommand{\hilfsmittel}{nicht programmierbarer Taschenrechner, 1 Formelblatt (A4)}

% ==========================================================================
%  PAKETE
% ==========================================================================
\usepackage[T1]{fontenc}
\usepackage[ngerman]{babel}
\usepackage[sc]{mathpazo}            % Mathe-Schrift passend zu Palatino
\usepackage{tgpagella}               % Fließtext: TeX Gyre Pagella
\usepackage[scale=0.92]{tgheros}     % Überschriften: TeX Gyre Heros
\renewcommand{\ttdefault}{lmtt}      % Code: Latin Modern Mono
\usepackage{microtype}
\usepackage{amsmath,amssymb}
\usepackage{xcolor}
\usepackage[most]{tcolorbox}
\usepackage{siunitx}
\sisetup{locale=DE, per-mode=symbol}
\usepackage{tikz}
\usetikzlibrary{arrows.meta,calc}
\usepackage{pgfplots}
\pgfplotsset{compat=1.18}
\usepackage[european]{circuitikz}
\usepackage{listings}
\usepackage[hidelinks]{hyperref}
\hypersetup{pdftitle={OhmTools Lernhelfer Übungsblatt für \modulname}}

% --- Farben ---------------------------------------------------------------
\definecolor{akzent}{HTML}{1F4E79}   % Hauptfarbe (dunkelblau)
\definecolor{akzenthell}{HTML}{E8EEF5}
\definecolor{grau}{HTML}{6B7280}
\definecolor{loesung}{HTML}{1B7F3B}
\definecolor{GridColor}{HTML}{C9D3DE}

% ==========================================================================
%  LAYOUT DER EXAM-KLASSE
% ==========================================================================
\extrawidth{0.4in}
\extraheadheight{-0.1in}
\extrafootheight{0.1in}

% Kopf- und Fußzeilen
\pagestyle{headandfoot}
\firstpageheader{}{}{}
\runningheader{\sffamily\small\color{grau}OhmTools Lernhelfer}
              {\ifprintanswers\sffamily\small\bfseries\color{loesung}LÖSUNG\fi}
              {\sffamily\small\color{grau}Übungsblatt \modulname}
\runningheadrule
\footer{}{\sffamily\small\color{grau}Seite \thepage\ von \numpages}{}

% Aufgabenkopf: „Aufgabe 1  Titel ............ 12 Punkte“
\newcommand{\aufgabenkopf}{%
  \begin{minipage}{\textwidth}
    \vspace{3mm}
    {\sffamily\bfseries\large\color{akzent}Aufgabe \thequestion\quad\thequestiontitle}%
    \hfill{\sffamily\small\color{grau}\totalpoints\ \points}\par\vspace{1mm}
    {\color{akzent}\rule{\textwidth}{0.8pt}}\rule[-3mm]{0pt}{0pt}
  \end{minipage}}
\qformat{\aufgabenkopf}

% Jede Aufgabe beginnt auf einer neuen Seite, die erste unter dem Titelkopf.
% So lässt sich jede Aufgabe am Tablet ohne Scrollen bearbeiten, der Rest
% der Seite bleibt frei.
\newcommand{\aufgabe}[1]{\ifnum\value{question}>0 \newpage\fi\titledquestion{#1}}

% Punkte bei Teilaufgaben
\pointpoints{Punkt}{Punkte}
\pointname{\,P.}
\pointformat{\textcolor{grau}{\sffamily\footnotesize[\thepoints]}}
\renewcommand{\partlabel}{\textbf{\thepartno)}}
\renewcommand{\partshook}{\setlength{\itemsep}{4pt}}

% Lösungen (erscheinen nur in der Lösungsversion)
\shadedsolutions
\definecolor{SolutionColor}{HTML}{EDF7F0}
\renewcommand{\solutiontitle}{\noindent\textbf{\sffamily\color{loesung}Lösung:}\enspace}

% Multiple Choice
\checkboxchar{$\square$}
\checkedchar{$\blacksquare$}
\CorrectChoiceEmphasis{\bfseries\color{loesung}}

% Code-Darstellung
\lstset{
  basicstyle=\ttfamily\small,
  columns=flexible, keepspaces=true,
  keywordstyle=\color{akzent}\bfseries,
  commentstyle=\color{grau}\itshape,
  stringstyle=\color{loesung},
  numbers=left, numberstyle=\tiny\color{grau}, numbersep=8pt,
  frame=leftline, rulecolor=\color{akzent}, framerule=1.2pt,
  xleftmargin=2.5em, aboveskip=6pt, belowskip=6pt,
  showstringspaces=false, tabsize=4,
  backgroundcolor=\color{akzenthell!60},
}

% Kasten für „Gegeben“
\newtcolorbox{gegeben}{enhanced, colback=akzenthell, colframe=akzenthell,
  boxrule=0pt, arc=2pt, left=3mm, right=3mm, top=1.5mm, bottom=1.5mm,
  fontupper=\small}

\begin{document}

% ==========================================================================
%  TITELKOPF
% ==========================================================================
\begin{tcolorbox}[enhanced, colback=akzent, colframe=akzent, coltext=white,
    arc=3pt, boxrule=0pt, left=6mm, right=6mm, top=4mm, bottom=4mm]
  {\sffamily\small OhmTools Lernhelfer\par}
  \vspace{2mm}
  {\sffamily\bfseries\LARGE Übungsblatt für \modulname\par}
  \ifprintanswers
    \vspace{2mm}
    {\sffamily\small\hfill\colorbox{white}{\sffamily\bfseries\color{loesung}\,LÖSUNGSVERSION\,}\par}
  \fi
\end{tcolorbox}

\vspace{1mm}
\noindent{\sffamily\small
  {\color{grau}Bearbeitungszeit:}~\dauer\hfill
  {\color{grau}Hilfsmittel:}~\hilfsmittel\hfill
  {\color{grau}Punkte:}~\numpoints\par}
\vspace{2mm}

% ==========================================================================
%  AUFGABEN
% ==========================================================================
\begin{questions}

% --------------------------------------------------------------------------
%  Aufgabe 1 – Mathematik
% --------------------------------------------------------------------------
\aufgabe{Kurvendiskussion}
Gegeben ist die Funktion
\[
  f(x) = x^3 - 6x^2 + 9x, \qquad x \in \mathbb{R}.
\]

\begin{parts}
\part[2] Bestimmen Sie die Nullstellen von $f$.
\begin{solution}
  $f(x) = x\,(x^2 - 6x + 9) = x\,(x-3)^2$.
  Nullstellen: $x_1 = 0$ (einfach), $x_2 = 3$ (doppelt).
\end{solution}

\part[4] Berechnen Sie die Extrempunkte und geben Sie jeweils deren Art an.
\begin{solution}
  $f'(x) = 3x^2 - 12x + 9 = 3\,(x-1)(x-3) \overset{!}{=} 0
   \;\Rightarrow\; x = 1$ oder $x = 3$.\\
  $f''(x) = 6x - 12$: \quad $f''(1) = -6 < 0 \Rightarrow$ Hochpunkt $H(1\,|\,4)$,
  \quad $f''(3) = 6 > 0 \Rightarrow$ Tiefpunkt $T(3\,|\,0)$.
\end{solution}

\part[2] Bestimmen Sie den Wendepunkt.
\begin{solution}
  $f''(x) = 0 \Rightarrow x = 2$, \ $f'''(2) = 6 \neq 0$.
  Wendepunkt $W(2\,|\,2)$.
\end{solution}

\part[3] Berechnen Sie den Inhalt der Fläche, die der Graph von $f$ im Intervall
$[0;\,3]$ mit der $x$-Achse einschließt.
\begin{solution}
  Da $f(x) = x(x-3)^2 \ge 0$ auf $[0;3]$:
  \[
    A = \int_0^3 f(x)\,\mathrm{d}x
      = \Big[\tfrac{x^4}{4} - 2x^3 + \tfrac{9}{2}x^2\Big]_0^3
      = \tfrac{81}{4} - 54 + \tfrac{81}{2} = \tfrac{27}{4} = 6{,}75 \text{ FE}.
  \]
\end{solution}

\part[2] Skizzieren Sie den Graphen von $f$ im folgenden Koordinatensystem
und markieren Sie die berechneten Punkte.

\begin{center}
\begin{tikzpicture}
  \begin{axis}[
      width=10cm, height=5.5cm,
      axis lines=middle, xmin=-0.6, xmax=4.6, ymin=-1.6, ymax=5.6,
      xtick={-0.5,0,...,4.5}, ytick={-1,...,5},
      xticklabels={,0,,1,,2,,3,,4,},
      grid=both, grid style={GridColor, line width=0.2pt},
      minor tick num=1, minor grid style={GridColor!50},
      xlabel={$x$}, ylabel={$y$},
      tick label style={font=\footnotesize},
      clip=true]
    \ifprintanswers
      \addplot[akzent, very thick, domain=-0.3:4.25, samples=120] {x^3 - 6*x^2 + 9*x};
      \addplot[only marks, mark=*, mark size=2pt, loesung]
        coordinates {(0,0) (3,0) (1,4) (2,2)};
      \node[loesung, above right, font=\footnotesize] at (axis cs:1,4) {$H$};
      \node[loesung, above right, font=\footnotesize] at (axis cs:2,2) {$W$};
      \node[loesung, below right, font=\footnotesize] at (axis cs:3,0) {$T$};
    \fi
  \end{axis}
\end{tikzpicture}
\end{center}
\end{parts}

% --------------------------------------------------------------------------
%  Aufgabe 2 – Elektrotechnik
% --------------------------------------------------------------------------
\aufgabe{Gleichstromnetzwerk}
Gegeben ist die folgende Schaltung.

\begin{center}
\begin{circuitikz}[line width=0.7pt]
  \draw (0,0) to[V, v<=$U_0$] (0,3)
        to[R, l=$R_1$, i>^=$I$] (3.5,3)
        to[short, -, i>^=$I_3$] (6.5,3)
        to[R, l_=$R_3$, v^=$U_P$] (6.5,0)
        -- (0,0);
  \draw (3.5,3) to[R, l=$R_2$, i>_=$I_2$] (3.5,0);
  \fill (3.5,3) circle (2pt) (3.5,0) circle (2pt);
\end{circuitikz}
\end{center}

\begin{gegeben}
  \textbf{Gegeben:}\quad
  $U_0 = \qty{12}{\volt}$, \quad
  $R_1 = \qty{100}{\ohm}$, \quad
  $R_2 = \qty{300}{\ohm}$, \quad
  $R_3 = \qty{150}{\ohm}$
\end{gegeben}

\begin{parts}
\part[3] Berechnen Sie den Gesamtwiderstand $R_\text{ges}$ der Schaltung.
\begin{solution}
  $R_{23} = \dfrac{R_2 R_3}{R_2 + R_3} = \dfrac{300 \cdot 150}{450}\,\unit{\ohm} = \qty{100}{\ohm}$,
  \quad $R_\text{ges} = R_1 + R_{23} = \qty{200}{\ohm}$.
\end{solution}

\part[2] Berechnen Sie den Gesamtstrom $I$.
\begin{solution}
  $I = \dfrac{U_0}{R_\text{ges}} = \dfrac{\qty{12}{\volt}}{\qty{200}{\ohm}} = \qty{60}{\milli\ampere}$.
\end{solution}

\part[4] Bestimmen Sie die Spannung $U_P$ an der Parallelschaltung sowie die
Teilströme $I_2$ und $I_3$.
\begin{solution}
  $U_P = I \cdot R_{23} = \qty{60}{\milli\ampere} \cdot \qty{100}{\ohm} = \qty{6}{\volt}$.\\
  $I_2 = U_P / R_2 = \qty{20}{\milli\ampere}$, \quad
  $I_3 = U_P / R_3 = \qty{40}{\milli\ampere}$
  \quad (Probe: $I_2 + I_3 = I$ \checkmark).
\end{solution}

\part[2] Welche Leistung wird im Widerstand $R_1$ umgesetzt?
\begin{solution}
  $P_1 = I^2 R_1 = (\qty{0,06}{\ampere})^2 \cdot \qty{100}{\ohm} = \qty{0,36}{\watt}$.
\end{solution}
\end{parts}

% --------------------------------------------------------------------------
%  Aufgabe 3 – Physik
% --------------------------------------------------------------------------
\aufgabe{Schräger Wurf}
Ein Ball wird vom Boden aus mit der Anfangsgeschwindigkeit $v_0$ unter dem Winkel
$\alpha$ gegen die Horizontale abgeworfen. Der Luftwiderstand wird vernachlässigt.

\begin{center}
\begin{tikzpicture}[>=Stealth]
  \draw[->, grau] (-0.3,0) -- (9.8,0) node[below] {$x$};
  \draw[->, grau] (0,-0.3) -- (0,3.4) node[left] {$y$};
  \draw[akzent, very thick] (0,0) parabola bend (4.5,2.6) (9,0);
  \draw[->, very thick] (0,0) -- (49:1.9) node[above left] {$\vec v_0$};
  \draw (1,0) arc (0:49:1);
  \node at (24:1.3) {$\alpha$};
  \draw[<->, dashed, grau] (4.5,0) -- node[right, black] {$h_\text{max}$} (4.5,2.6);
  \draw[<->] (0,-0.6) -- node[fill=white, inner sep=2pt] {$w$} (9,-0.6);
  \fill[akzent] (4.5,2.6) circle (2pt);
  \node[grau, font=\footnotesize, anchor=north east] at (9.8,3.4) {Skizze nicht maßstäblich};
\end{tikzpicture}
\end{center}

\begin{gegeben}
  \textbf{Gegeben:}\quad
  $v_0 = \qty{20}{\metre\per\second}$, \quad
  $\alpha = \ang{30}$, \quad
  $g = \qty{9,81}{\metre\per\second\squared}$
\end{gegeben}

\begin{parts}
\part[2] Nach welcher Zeit $t_s$ erreicht der Ball seine maximale Höhe?
\begin{solution}
  Am Scheitel gilt $v_y = 0$:\quad
  $t_s = \dfrac{v_0 \sin\alpha}{g} = \dfrac{\qty{10}{\metre\per\second}}{\qty{9,81}{\metre\per\second\squared}} \approx \qty{1,02}{\second}$.
\end{solution}

\part[3] Berechnen Sie die maximale Höhe $h_\text{max}$.
\begin{solution}
  $h_\text{max} = \dfrac{v_0^2 \sin^2\alpha}{2g}
   = \dfrac{(\qty{10}{\metre\per\second})^2}{2 \cdot \qty{9,81}{\metre\per\second\squared}}
   \approx \qty{5,10}{\metre}$.
\end{solution}

\part[3] Berechnen Sie die Wurfweite $w$.
\begin{solution}
  $w = \dfrac{v_0^2 \sin(2\alpha)}{g}
   = \dfrac{\qty{400}{\metre\squared\per\second\squared} \cdot \sin\ang{60}}{\qty{9,81}{\metre\per\second\squared}}
   \approx \qty{35,3}{\metre}$.
\end{solution}

\part[2] Unter welchem Abwurfwinkel wird die Wurfweite maximal? Begründen Sie kurz.
\begin{solution}
  $w \sim \sin(2\alpha)$ wird maximal für $2\alpha = \ang{90}$, also $\alpha = \ang{45}$.
\end{solution}
\end{parts}

% --------------------------------------------------------------------------
%  Aufgabe 4 – Informatik
% --------------------------------------------------------------------------
\aufgabe{Rekursion und Zahlensysteme}
Betrachten Sie die folgende Python-Funktion:

\begin{lstlisting}[language=Python]
def f(n):
    if n <= 1:
        return n
    return f(n - 1) + f(n - 2)
\end{lstlisting}

\begin{parts}
\part[2] Welchen Wert liefert \lstinline{f(6)}? Geben Sie die Zwischenwerte an.
\begin{solution}
  $f(0..6) = 0,\,1,\,1,\,2,\,3,\,5,\,8$, \ also \lstinline{f(6)} $= 8$.
\end{solution}

\part[1] Was berechnet die Funktion allgemein?
\begin{solution}
  Die $n$-te Fibonacci-Zahl.
\end{solution}

\part[3] Geben Sie die Laufzeitkomplexität in $\mathcal{O}$-Notation an und
begründen Sie. Wie lässt sich die Laufzeit verbessern?
\begin{solution}
  $\mathcal{O}(2^n)$: Jeder Aufruf erzeugt bis zu zwei weitere Aufrufe, der
  Rekursionsbaum hat die Tiefe $n$, und Teilergebnisse werden mehrfach berechnet.
  Mit Memoisation oder einer iterativen Lösung sinkt die Laufzeit auf $\mathcal{O}(n)$.
\end{solution}

\part[1] Welche Datenstruktur arbeitet nach dem LIFO-Prinzip? Kreuzen Sie an.
\par\smallskip
\begin{oneparcheckboxes}
  \choice Warteschlange (Queue)
  \CorrectChoice Stapel (Stack)
  \choice Heap
  \choice Hashtabelle
\end{oneparcheckboxes}

\part[2] Wandeln Sie die Dezimalzahl $45$ in das Binär- und das Hexadezimalsystem um.
\begin{solution}
  $45 = 32 + 8 + 4 + 1 = 101101_2$, \quad $0010\,1101_2 = 2\mathrm{D}_{16}$.
\end{solution}
\end{parts}

\end{questions}

\end{document}
"
%                    (ebenfalls 3x ausführen)   – oder einfach: sh build.sh
%
%  Alternativ unten die Zeile  % \printanswers  einkommentieren.
%
%  Benötigt: TeX Live mit exam, circuitikz, pgfplots, siunitx, tcolorbox,
%            listings, tex-gyre (Pagella/Heros), lmodern (Code-Schrift),
%            babel-german, hyperref
% ==========================================================================
\documentclass[11pt,a4paper,addpoints]{exam}

% --- Lösungen ein/aus ----------------------------------------------------
% \printanswers
\ifdefined\mitloesung\printanswers\fi

% ==========================================================================
%  KOPFDATEN – hier anpassen
% ==========================================================================
\newcommand{\modulname}{Ingenieurgrundlagen}
\newcommand{\dauer}{90 Minuten}
\newcommand{\hilfsmittel}{nicht programmierbarer Taschenrechner, 1 Formelblatt (A4)}

% ==========================================================================
%  PAKETE
% ==========================================================================
\usepackage[T1]{fontenc}
\usepackage[ngerman]{babel}
\usepackage[sc]{mathpazo}            % Mathe-Schrift passend zu Palatino
\usepackage{tgpagella}               % Fließtext: TeX Gyre Pagella
\usepackage[scale=0.92]{tgheros}     % Überschriften: TeX Gyre Heros
\renewcommand{\ttdefault}{lmtt}      % Code: Latin Modern Mono
\usepackage{microtype}
\usepackage{amsmath,amssymb}
\usepackage{xcolor}
\usepackage[most]{tcolorbox}
\usepackage{siunitx}
\sisetup{locale=DE, per-mode=symbol}
\usepackage{tikz}
\usetikzlibrary{arrows.meta,calc}
\usepackage{pgfplots}
\pgfplotsset{compat=1.18}
\usepackage[european]{circuitikz}
\usepackage{listings}
\usepackage[hidelinks]{hyperref}
\hypersetup{pdftitle={OhmTools Lernhelfer Übungsblatt für \modulname}}

% --- Farben ---------------------------------------------------------------
\definecolor{akzent}{HTML}{1F4E79}   % Hauptfarbe (dunkelblau)
\definecolor{akzenthell}{HTML}{E8EEF5}
\definecolor{grau}{HTML}{6B7280}
\definecolor{loesung}{HTML}{1B7F3B}
\definecolor{GridColor}{HTML}{C9D3DE}

% ==========================================================================
%  LAYOUT DER EXAM-KLASSE
% ==========================================================================
\extrawidth{0.4in}
\extraheadheight{-0.1in}
\extrafootheight{0.1in}

% Kopf- und Fußzeilen
\pagestyle{headandfoot}
\firstpageheader{}{}{}
\runningheader{\sffamily\small\color{grau}OhmTools Lernhelfer}
              {\ifprintanswers\sffamily\small\bfseries\color{loesung}LÖSUNG\fi}
              {\sffamily\small\color{grau}Übungsblatt \modulname}
\runningheadrule
\footer{}{\sffamily\small\color{grau}Seite \thepage\ von \numpages}{}

% Aufgabenkopf: „Aufgabe 1  Titel ............ 12 Punkte“
\newcommand{\aufgabenkopf}{%
  \begin{minipage}{\textwidth}
    \vspace{3mm}
    {\sffamily\bfseries\large\color{akzent}Aufgabe \thequestion\quad\thequestiontitle}%
    \hfill{\sffamily\small\color{grau}\totalpoints\ \points}\par\vspace{1mm}
    {\color{akzent}\rule{\textwidth}{0.8pt}}\rule[-3mm]{0pt}{0pt}
  \end{minipage}}
\qformat{\aufgabenkopf}

% Jede Aufgabe beginnt auf einer neuen Seite, die erste unter dem Titelkopf.
% So lässt sich jede Aufgabe am Tablet ohne Scrollen bearbeiten, der Rest
% der Seite bleibt frei.
\newcommand{\aufgabe}[1]{\ifnum\value{question}>0 \newpage\fi\titledquestion{#1}}

% Punkte bei Teilaufgaben
\pointpoints{Punkt}{Punkte}
\pointname{\,P.}
\pointformat{\textcolor{grau}{\sffamily\footnotesize[\thepoints]}}
\renewcommand{\partlabel}{\textbf{\thepartno)}}
\renewcommand{\partshook}{\setlength{\itemsep}{4pt}}

% Lösungen (erscheinen nur in der Lösungsversion)
\shadedsolutions
\definecolor{SolutionColor}{HTML}{EDF7F0}
\renewcommand{\solutiontitle}{\noindent\textbf{\sffamily\color{loesung}Lösung:}\enspace}

% Multiple Choice
\checkboxchar{$\square$}
\checkedchar{$\blacksquare$}
\CorrectChoiceEmphasis{\bfseries\color{loesung}}

% Code-Darstellung
\lstset{
  basicstyle=\ttfamily\small,
  columns=flexible, keepspaces=true,
  keywordstyle=\color{akzent}\bfseries,
  commentstyle=\color{grau}\itshape,
  stringstyle=\color{loesung},
  numbers=left, numberstyle=\tiny\color{grau}, numbersep=8pt,
  frame=leftline, rulecolor=\color{akzent}, framerule=1.2pt,
  xleftmargin=2.5em, aboveskip=6pt, belowskip=6pt,
  showstringspaces=false, tabsize=4,
  backgroundcolor=\color{akzenthell!60},
}

% Kasten für „Gegeben“
\newtcolorbox{gegeben}{enhanced, colback=akzenthell, colframe=akzenthell,
  boxrule=0pt, arc=2pt, left=3mm, right=3mm, top=1.5mm, bottom=1.5mm,
  fontupper=\small}

\begin{document}

% ==========================================================================
%  TITELKOPF
% ==========================================================================
\begin{tcolorbox}[enhanced, colback=akzent, colframe=akzent, coltext=white,
    arc=3pt, boxrule=0pt, left=6mm, right=6mm, top=4mm, bottom=4mm]
  {\sffamily\small OhmTools Lernhelfer\par}
  \vspace{2mm}
  {\sffamily\bfseries\LARGE Übungsblatt für \modulname\par}
  \ifprintanswers
    \vspace{2mm}
    {\sffamily\small\hfill\colorbox{white}{\sffamily\bfseries\color{loesung}\,LÖSUNGSVERSION\,}\par}
  \fi
\end{tcolorbox}

\vspace{1mm}
\noindent{\sffamily\small
  {\color{grau}Bearbeitungszeit:}~\dauer\hfill
  {\color{grau}Hilfsmittel:}~\hilfsmittel\hfill
  {\color{grau}Punkte:}~\numpoints\par}
\vspace{2mm}

% ==========================================================================
%  AUFGABEN
% ==========================================================================
\begin{questions}

% --------------------------------------------------------------------------
%  Aufgabe 1 – Mathematik
% --------------------------------------------------------------------------
\aufgabe{Kurvendiskussion}
Gegeben ist die Funktion
\[
  f(x) = x^3 - 6x^2 + 9x, \qquad x \in \mathbb{R}.
\]

\begin{parts}
\part[2] Bestimmen Sie die Nullstellen von $f$.
\begin{solution}
  $f(x) = x\,(x^2 - 6x + 9) = x\,(x-3)^2$.
  Nullstellen: $x_1 = 0$ (einfach), $x_2 = 3$ (doppelt).
\end{solution}

\part[4] Berechnen Sie die Extrempunkte und geben Sie jeweils deren Art an.
\begin{solution}
  $f'(x) = 3x^2 - 12x + 9 = 3\,(x-1)(x-3) \overset{!}{=} 0
   \;\Rightarrow\; x = 1$ oder $x = 3$.\\
  $f''(x) = 6x - 12$: \quad $f''(1) = -6 < 0 \Rightarrow$ Hochpunkt $H(1\,|\,4)$,
  \quad $f''(3) = 6 > 0 \Rightarrow$ Tiefpunkt $T(3\,|\,0)$.
\end{solution}

\part[2] Bestimmen Sie den Wendepunkt.
\begin{solution}
  $f''(x) = 0 \Rightarrow x = 2$, \ $f'''(2) = 6 \neq 0$.
  Wendepunkt $W(2\,|\,2)$.
\end{solution}

\part[3] Berechnen Sie den Inhalt der Fläche, die der Graph von $f$ im Intervall
$[0;\,3]$ mit der $x$-Achse einschließt.
\begin{solution}
  Da $f(x) = x(x-3)^2 \ge 0$ auf $[0;3]$:
  \[
    A = \int_0^3 f(x)\,\mathrm{d}x
      = \Big[\tfrac{x^4}{4} - 2x^3 + \tfrac{9}{2}x^2\Big]_0^3
      = \tfrac{81}{4} - 54 + \tfrac{81}{2} = \tfrac{27}{4} = 6{,}75 \text{ FE}.
  \]
\end{solution}

\part[2] Skizzieren Sie den Graphen von $f$ im folgenden Koordinatensystem
und markieren Sie die berechneten Punkte.

\begin{center}
\begin{tikzpicture}
  \begin{axis}[
      width=10cm, height=5.5cm,
      axis lines=middle, xmin=-0.6, xmax=4.6, ymin=-1.6, ymax=5.6,
      xtick={-0.5,0,...,4.5}, ytick={-1,...,5},
      xticklabels={,0,,1,,2,,3,,4,},
      grid=both, grid style={GridColor, line width=0.2pt},
      minor tick num=1, minor grid style={GridColor!50},
      xlabel={$x$}, ylabel={$y$},
      tick label style={font=\footnotesize},
      clip=true]
    \ifprintanswers
      \addplot[akzent, very thick, domain=-0.3:4.25, samples=120] {x^3 - 6*x^2 + 9*x};
      \addplot[only marks, mark=*, mark size=2pt, loesung]
        coordinates {(0,0) (3,0) (1,4) (2,2)};
      \node[loesung, above right, font=\footnotesize] at (axis cs:1,4) {$H$};
      \node[loesung, above right, font=\footnotesize] at (axis cs:2,2) {$W$};
      \node[loesung, below right, font=\footnotesize] at (axis cs:3,0) {$T$};
    \fi
  \end{axis}
\end{tikzpicture}
\end{center}
\end{parts}

% --------------------------------------------------------------------------
%  Aufgabe 2 – Elektrotechnik
% --------------------------------------------------------------------------
\aufgabe{Gleichstromnetzwerk}
Gegeben ist die folgende Schaltung.

\begin{center}
\begin{circuitikz}[line width=0.7pt]
  \draw (0,0) to[V, v<=$U_0$] (0,3)
        to[R, l=$R_1$, i>^=$I$] (3.5,3)
        to[short, -, i>^=$I_3$] (6.5,3)
        to[R, l_=$R_3$, v^=$U_P$] (6.5,0)
        -- (0,0);
  \draw (3.5,3) to[R, l=$R_2$, i>_=$I_2$] (3.5,0);
  \fill (3.5,3) circle (2pt) (3.5,0) circle (2pt);
\end{circuitikz}
\end{center}

\begin{gegeben}
  \textbf{Gegeben:}\quad
  $U_0 = \qty{12}{\volt}$, \quad
  $R_1 = \qty{100}{\ohm}$, \quad
  $R_2 = \qty{300}{\ohm}$, \quad
  $R_3 = \qty{150}{\ohm}$
\end{gegeben}

\begin{parts}
\part[3] Berechnen Sie den Gesamtwiderstand $R_\text{ges}$ der Schaltung.
\begin{solution}
  $R_{23} = \dfrac{R_2 R_3}{R_2 + R_3} = \dfrac{300 \cdot 150}{450}\,\unit{\ohm} = \qty{100}{\ohm}$,
  \quad $R_\text{ges} = R_1 + R_{23} = \qty{200}{\ohm}$.
\end{solution}

\part[2] Berechnen Sie den Gesamtstrom $I$.
\begin{solution}
  $I = \dfrac{U_0}{R_\text{ges}} = \dfrac{\qty{12}{\volt}}{\qty{200}{\ohm}} = \qty{60}{\milli\ampere}$.
\end{solution}

\part[4] Bestimmen Sie die Spannung $U_P$ an der Parallelschaltung sowie die
Teilströme $I_2$ und $I_3$.
\begin{solution}
  $U_P = I \cdot R_{23} = \qty{60}{\milli\ampere} \cdot \qty{100}{\ohm} = \qty{6}{\volt}$.\\
  $I_2 = U_P / R_2 = \qty{20}{\milli\ampere}$, \quad
  $I_3 = U_P / R_3 = \qty{40}{\milli\ampere}$
  \quad (Probe: $I_2 + I_3 = I$ \checkmark).
\end{solution}

\part[2] Welche Leistung wird im Widerstand $R_1$ umgesetzt?
\begin{solution}
  $P_1 = I^2 R_1 = (\qty{0,06}{\ampere})^2 \cdot \qty{100}{\ohm} = \qty{0,36}{\watt}$.
\end{solution}
\end{parts}

% --------------------------------------------------------------------------
%  Aufgabe 3 – Physik
% --------------------------------------------------------------------------
\aufgabe{Schräger Wurf}
Ein Ball wird vom Boden aus mit der Anfangsgeschwindigkeit $v_0$ unter dem Winkel
$\alpha$ gegen die Horizontale abgeworfen. Der Luftwiderstand wird vernachlässigt.

\begin{center}
\begin{tikzpicture}[>=Stealth]
  \draw[->, grau] (-0.3,0) -- (9.8,0) node[below] {$x$};
  \draw[->, grau] (0,-0.3) -- (0,3.4) node[left] {$y$};
  \draw[akzent, very thick] (0,0) parabola bend (4.5,2.6) (9,0);
  \draw[->, very thick] (0,0) -- (49:1.9) node[above left] {$\vec v_0$};
  \draw (1,0) arc (0:49:1);
  \node at (24:1.3) {$\alpha$};
  \draw[<->, dashed, grau] (4.5,0) -- node[right, black] {$h_\text{max}$} (4.5,2.6);
  \draw[<->] (0,-0.6) -- node[fill=white, inner sep=2pt] {$w$} (9,-0.6);
  \fill[akzent] (4.5,2.6) circle (2pt);
  \node[grau, font=\footnotesize, anchor=north east] at (9.8,3.4) {Skizze nicht maßstäblich};
\end{tikzpicture}
\end{center}

\begin{gegeben}
  \textbf{Gegeben:}\quad
  $v_0 = \qty{20}{\metre\per\second}$, \quad
  $\alpha = \ang{30}$, \quad
  $g = \qty{9,81}{\metre\per\second\squared}$
\end{gegeben}

\begin{parts}
\part[2] Nach welcher Zeit $t_s$ erreicht der Ball seine maximale Höhe?
\begin{solution}
  Am Scheitel gilt $v_y = 0$:\quad
  $t_s = \dfrac{v_0 \sin\alpha}{g} = \dfrac{\qty{10}{\metre\per\second}}{\qty{9,81}{\metre\per\second\squared}} \approx \qty{1,02}{\second}$.
\end{solution}

\part[3] Berechnen Sie die maximale Höhe $h_\text{max}$.
\begin{solution}
  $h_\text{max} = \dfrac{v_0^2 \sin^2\alpha}{2g}
   = \dfrac{(\qty{10}{\metre\per\second})^2}{2 \cdot \qty{9,81}{\metre\per\second\squared}}
   \approx \qty{5,10}{\metre}$.
\end{solution}

\part[3] Berechnen Sie die Wurfweite $w$.
\begin{solution}
  $w = \dfrac{v_0^2 \sin(2\alpha)}{g}
   = \dfrac{\qty{400}{\metre\squared\per\second\squared} \cdot \sin\ang{60}}{\qty{9,81}{\metre\per\second\squared}}
   \approx \qty{35,3}{\metre}$.
\end{solution}

\part[2] Unter welchem Abwurfwinkel wird die Wurfweite maximal? Begründen Sie kurz.
\begin{solution}
  $w \sim \sin(2\alpha)$ wird maximal für $2\alpha = \ang{90}$, also $\alpha = \ang{45}$.
\end{solution}
\end{parts}

% --------------------------------------------------------------------------
%  Aufgabe 4 – Informatik
% --------------------------------------------------------------------------
\aufgabe{Rekursion und Zahlensysteme}
Betrachten Sie die folgende Python-Funktion:

\begin{lstlisting}[language=Python]
def f(n):
    if n <= 1:
        return n
    return f(n - 1) + f(n - 2)
\end{lstlisting}

\begin{parts}
\part[2] Welchen Wert liefert \lstinline{f(6)}? Geben Sie die Zwischenwerte an.
\begin{solution}
  $f(0..6) = 0,\,1,\,1,\,2,\,3,\,5,\,8$, \ also \lstinline{f(6)} $= 8$.
\end{solution}

\part[1] Was berechnet die Funktion allgemein?
\begin{solution}
  Die $n$-te Fibonacci-Zahl.
\end{solution}

\part[3] Geben Sie die Laufzeitkomplexität in $\mathcal{O}$-Notation an und
begründen Sie. Wie lässt sich die Laufzeit verbessern?
\begin{solution}
  $\mathcal{O}(2^n)$: Jeder Aufruf erzeugt bis zu zwei weitere Aufrufe, der
  Rekursionsbaum hat die Tiefe $n$, und Teilergebnisse werden mehrfach berechnet.
  Mit Memoisation oder einer iterativen Lösung sinkt die Laufzeit auf $\mathcal{O}(n)$.
\end{solution}

\part[1] Welche Datenstruktur arbeitet nach dem LIFO-Prinzip? Kreuzen Sie an.
\par\smallskip
\begin{oneparcheckboxes}
  \choice Warteschlange (Queue)
  \CorrectChoice Stapel (Stack)
  \choice Heap
  \choice Hashtabelle
\end{oneparcheckboxes}

\part[2] Wandeln Sie die Dezimalzahl $45$ in das Binär- und das Hexadezimalsystem um.
\begin{solution}
  $45 = 32 + 8 + 4 + 1 = 101101_2$, \quad $0010\,1101_2 = 2\mathrm{D}_{16}$.
\end{solution}
\end{parts}

\end{questions}

\end{document}
